\documentclass[12pt]{article}
\usepackage[T1]{fontenc}
\usepackage[utf8]{inputenc}
\usepackage{lmodern}
\usepackage{textcomp}
\usepackage{enumitem}
\usepackage{geometry}
\usepackage{graphicx}
\usepackage{amsmath, amssymb}
\geometry{margin=1in}
\begin{document}
\begin{center}
Astronomy - Graduate Level Assessment 23 \
Total Marks: 40 \
Answer all questions.
\end{center}

\section*{Multiple Choice (10 marks)}
\begin{enumerate}[label=\arabic*.]
\item According to the Solar Nebular theory what are asteroids and comets?
\begin{enumerate}[label=(\alph*)]
    \item They are the shattered remains of collisions between planets.
    \item They are chunks of rock or ice that condensed long after the planets and moons had formed.
    \item They are chunks of rock or ice that were expelled from planets by volcanoes.
    \item They are leftover planetesimals that never accreted into planets.
\end{enumerate}
\item Approximately how far away is the Andromeda Galaxy?
\begin{enumerate}[label=(\alph*)]
    \item 1.7 million light years
    \item 2.1 million light years
    \item 2.5 million light years
    \item 3.2 million light years
\end{enumerate}
\item Moons cause/contribute to which of the following?
\begin{enumerate}[label=(\alph*)]
    \item stability of particles within rings.
    \item gravitational effects at ring edges as the moons pass by.
    \item gaps between rings.
    \item Moons contribute to all of the above.
\end{enumerate}
\item What is not true of Jupiter's magnetic field?
\begin{enumerate}[label=(\alph*)]
    \item it is about 20000 times stronger than Earth's magnetic field
    \item it traps charged particles from Io's volcanoes in a "plasma torus" around the planet
    \item it causes spectacular auroral displays at Jupiter's north and south poles
    \item it switches polarity every 11 years
\end{enumerate}
\item How do astronomers think Jupiter generates its internal heat?
\begin{enumerate}[label=(\alph*)]
    \item nuclear fusion in the core
    \item by contracting changing gravitational potential energy into thermal energy
    \item internal friction due to its high rotation rate
    \item chemical processes
\end{enumerate}
\end{enumerate}

\section*{True / False (10 marks)}
\textit{Answer True or False}
\begin{enumerate}[label=\arabic*.]
\setcounter{enumi}{5}
\item True or False: The nebular theory does not predict an equal number of terrestrial and jovian planets in the solar system.
\item True or False: Volcanic heating cannot account for Mercury's lack of a permanent atmosphere due to its insufficient volcanic activity.
\item True or False: A solar day on Planet X lasts approximately 100 Earth days, equal to its sidereal day length.
\item True or False: Type-Ia supernovae typically produce high amounts of X-rays, making them prominent sources of X-ray emissions in the universe.
\item True or False: The azimuth of the north celestial pole (NCP) indicates your geographic latitude in the sky.
\end{enumerate}

\section*{Long Form Response (20 marks)}
\textit{Answer each question with a clear and complete explanation.}
\begin{enumerate}[label=\arabic*.]
\setcounter{enumi}{10}
\item Discuss the unique characteristics of asteroids that differentiate them from background stars in sky surveys, focusing on their motion and implications for observation techniques.
\item Discuss the formation of "hot Jupiters" around other stars, analyzing their origin beyond the frost line and the mechanisms that lead to their inward migration.
\end{enumerate}

\vfill
\noindent\textit{End of Paper}
\end{document}